% problem_id: S001-lemma-non-effective-descent-projective
% source_id: S001 (Stacks Project)
% source_file: examples.tex
% statement_locator: examples.tex lines 5241--5247
% proof_locator: examples.tex lines 5249--5369
% env: lemma
% id_note: label 在本来源内唯一
% pairing: tight (verified)
% status: complete (statement + author proof verbatim + reason + locators)
%
\begin{lemma}
\label{lemma-non-effective-descent-projective}
There is an etale covering $X\to S$ of schemes and a descent datum
$(V/X,\varphi)$ relative to $X\to S$ such that 
$V\to X$ is projective,
but the descent datum is not effective in the category of schemes.
\end{lemma}

\begin{proof}
We imitate Hironaka's example of a smooth separated complex
algebraic space of dimension 3
which is not a scheme \cite[Example B.3.4.2]{H}.

\medskip\noindent
Consider the action of the group $G = \mathbf{Z}/2 = \{1, g\}$
on projective 3-space
$\mathbf{P}^3$ over the complex numbers by
$$
g[x,y,z,w] = [y,x,w,z].
$$
The action is free outside the two disjoint lines
$L_1=\{ [x,x,z,z]\}$ and $L_2=\{ [x,-x,z,-z]\}$ in
${\mathbf P}^3$. Let $Y={\mathbf P}^3-(L_1\cup L_2)$. There is a
smooth quasi-projective scheme $S=Y/G$ over ${\mathbf C}$ such that
$Y\to S$ is a $G$-torsor (Groupoids,
Definition \ref{groupoids-definition-principal-homogeneous-space}).
Explicitly, we can define $S$ as the image of the open subset $Y$
in ${\mathbf P}^3$ under the morphism
\begin{align*}
{\mathbf P}^3 & \to \text{Proj } {\mathbf C}[x,y,z,w]^G\\
   & = \text{Proj } {\mathbf C}[u_0,u_1,v_0,v_1,v_2]/(v_0v_1=v_2^2),
\end{align*}
where $u_0=x+y$, $u_1=z+w$, $v_0=(x-y)^2$, $v_1=(z-w)^2$,
and $v_2=(x-y)(z-w)$, and the ring is graded with $u_0,u_1$
in degree 1 and $v_0,v_1,v_2$ in degree 2.

\medskip\noindent
Let $C=\{ [x,y,z,w]: xy=z^2, w=0\}$ and $D=\{ [x,y,z,w]:
xy=w^2, z=0\}$. These are smooth conic curves in ${\mathbf P}^3$, contained
in the $G$-invariant open subset $Y$, with $g(C)=D$. Also,
$C\cap D$ consists of the two points $P:=[1,0,0,0]$
and $Q:=[0,1,0,0]$, and these two points are switched by the action
of $G$. 

\medskip\noindent
Let $V_Y\to Y$ be the scheme which over $Y-P$
is defined by blowing up $D$ and then the strict transform
of $C$, and over $Y-Q$ is defined by blowing up $C$ and then
the strict transform of $D$. (This is the same construction
as in the proof of Lemma \ref{lemma-non-descending-property-projective},
except that $Y$ here denotes an open subset of ${\mathbf P}^3$
rather than all of ${\mathbf P}^3$.)
Then the action of $G$ on $Y$ lifts to an action of $G$ on $V_Y$,
which switches the inverse images of $Y-P$ and $Y-Q$. This action
of $G$ on $V_Y$ gives a descent datum $(V_Y/Y,\varphi_Y)$
on $V_Y$ relative to the $G$-torsor
$Y\to S$. The morphism $V_Y\to Y$
is proper but not projective, as shown in the proof
of Lemma \ref{lemma-non-descending-property-projective}.

\medskip\noindent
Let $X$ be the disjoint union of the open subsets $Y-P$ and $Y-Q$;
then we have surjective etale morphisms $X\to Y\to S$.
Let $V$ be the pullback of $V_Y\to Y$ to $X$; then the morphism
$V\to X$ is projective, since $V_Y\to Y$ is a blowup over each of the open
subsets $Y-P$ and $Y-Q$. Moreover, the descent datum $(V_Y/Y,\varphi_Y)$
pulls back to a descent datum $(V/X,\varphi)$ relative to the
etale covering $X\to S$.

\medskip\noindent
Suppose that this descent datum is effective in the category
of schemes. That is, there is a scheme $U\to S$
which pulls back to the morphism $V\to X$ together
with its descent datum. Then $U$ would be the quotient
of $V_Y$ by its $G$-action.
$$
\xymatrix{
V \ar[r]\ar[d]&  X\ar[d] \\
V_Y \ar[r]\ar[d]& Y\ar[d] \\
U \ar[r]& S
}
$$

\medskip\noindent
Let $E$ be the inverse image of $C\cup D\subset Y$ in $V_Y$;
thus $E\rightarrow C\cup D$ is a proper morphism, with fibers
isomorphic to ${\mathbf P}^1$ over $(C\cup D)-\{P,Q\}$.
The inverse image of $P$ in $E$ is a union of two lines $L_0$
and $M_0$. It follows that the inverse image of $Q=g(P)$ in $E$
is the union of two lines $L_0'=g(M_0)$ and $M_0'=g(L_0)$.
As shown in the proof
of Lemma \ref{lemma-non-descending-property-projective},
we have a rational equivalence $L_0+M_0'=L_0+g(L_0)\sim 0$ on $E$.

\medskip\noindent
By descent of closed subschemes, there is a curve $L_1\subset U$
(isomorphic to ${\mathbf P}^1$)
whose inverse image in $V_Y$ is $L_0\cup g(L_0)$. (Use Descent, Lemma
\ref{descent-lemma-affine}, noting that a closed immersion is an affine
morphism.)
Let $R$ be a complex point of $L_1$. Since
we assumed that $U$ is a scheme, we can choose a function
$f$ in the local ring $O_{U,R}$ that vanishes at $R$ but not
on the whole curve $L_1$. Let $D_{\text{loc}}$ be an irreducible component
of the closed subset $\{f = 0\}$ in $\Spec O_{U,R}$; then
$D_{\text{loc}}$ has codimension 1.
The closure of $D_{\text{loc}}$ in $U$ is an irreducible divisor $D_U$
in $U$ which contains the point $R$ but not the whole curve $L_1$.
The inverse image of $D_U$ in $V_Y$ is an effective divisor $D$
which intersects $L_0\cup g(L_0)$ but does not contain either
curve $L_0$ or $g(L_0)$.

\medskip\noindent
Since the complex 3-fold $V_Y$ is smooth, $O(D)$ is a line
bundle on $V_Y$. We use here that a regular local ring is factorial,
or in other words is a UFD, see
More on Algebra, Lemma \ref{more-algebra-lemma-regular-local-UFD}.
The restriction of $O(D)$ to the proper surface
$E\subset V_Y$ is a line bundle which has positive degree on the 1-cycle
$L_0+g(L_0)$, by our information on $D$. Since
$L_0+g(L_0)\sim 0$ on $E$, this contradicts 
that the degree of a line bundle is well-defined
on 1-cycles modulo rational equivalence on a proper scheme
over a field (Chow Homology,
Lemma \ref{chow-lemma-proper-pushforward-rational-equivalence}
and Lemma \ref{chow-lemma-factors}). Therefore the descent datum
$(V/X,\varphi)$ is in fact not effective; that is, $U$ does not exist
as a scheme.
\end{proof}
